% GENERATED FILE. Edit canonical lessons, then run npm run generate:books.
% canonical-source-hash: fnv1a64:e939144a35c193f3
% canonical-lessons: RU-C158-zhivoy, RU-C158-myortvyy, RU-C158-golodnyy, RU-C158-sytyy, RU-C158-znakomyy

\chapter{Alive, Dead, Hungry, Well-fed, Familiar}
\label{ch:ru-alive-dead-hungry-well-fed-familiar}
\hlchaptermodality{158}

\begin{chapteropening}
\textbf{By the end of this chapter:} \emph{I can say alive, dead, hungry, well-fed and familiar.}

Complete the last lesson of chapter 158: I can say alive, dead, hungry, well-fed and familiar.
\end{chapteropening}

\section[\texorpdfstring{\ru{живой} (zhivóy)}{живой (zhivóy)}]{\ru{живой} (zhivóy) — alive, lively}

\label{lesson:RU-C158-zhivoy}



\begin{quote}
Before the new one: say the Russian for tasty, then the Russian for comfortable, convenient.
\end{quote}

\subsection*{You'll want to know: \ru{живой}}
\textbf{\ru{живой}} (\emph{zhivóy}) — \textquotedblleft{}alive, lively\textquotedblright{}.

\begin{grammarlens}[title={What you've built}]
One more word for this chapter.
\end{grammarlens}

\subsection*{Your turn}
\begin{itemize}
\raggedright
  \item \emph{Say it:} \emph{zhivóy}
  \item \emph{Say it:} \emph{zhivóy}, once more
  \item \emph{Say it:} say \emph{zhivóy}
  \item \emph{Recall:} say the Russian for tasty, then the Russian for comfortable, convenient, then say \emph{zhivóy} again
\end{itemize}

\begin{tcolorbox}[breakable,colback=teal!4,colframe=teal!35!black,title={Before you move on}]
What does \ru{живой} mean? (\textquotedblleft{}Alive, lively\textquotedblright{}.) Say it once more.
\end{tcolorbox}
\section[\texorpdfstring{\ru{мёртвый} (myórtvyy)}{мёртвый (myórtvyy)}]{\ru{мёртвый} (myórtvyy) — dead}

\label{lesson:RU-C158-myortvyy}



\begin{quote}
Before the new one: say the Russian for comfortable, convenient, then the Russian for alive.
\end{quote}

\subsection*{You'll want to know: \ru{мёртвый}}
\textbf{\ru{мёртвый}} (\emph{myórtvyy}) — \textquotedblleft{}dead\textquotedblright{}.

\begin{grammarlens}[title={What you've built}]
One more word for this chapter.
\end{grammarlens}

\subsection*{Your turn}
\begin{itemize}
\raggedright
  \item \emph{Say it:} \emph{myórtvyy}
  \item \emph{Say it:} \emph{myórtvyy}, once more
  \item \emph{Say it:} say \emph{myórtvyy}
  \item \emph{Recall:} say the Russian for comfortable, convenient, then the Russian for alive, then say \emph{myórtvyy} again
\end{itemize}

\begin{tcolorbox}[breakable,colback=teal!4,colframe=teal!35!black,title={Before you move on}]
What does \ru{мёртвый} mean? (\textquotedblleft{}Dead\textquotedblright{}.) Say it once more.
\end{tcolorbox}
\section[\texorpdfstring{\ru{голодный} (golódnyy)}{голодный (golódnyy)}]{\ru{голодный} (golódnyy) — hungry}

\label{lesson:RU-C158-golodnyy}



\begin{quote}
Before the new one: say the Russian for alive, then the Russian for dead.
\end{quote}

\subsection*{You'll want to know: \ru{голодный}}
\textbf{\ru{голодный}} (\emph{golódnyy}) — \textquotedblleft{}hungry\textquotedblright{}.

\begin{grammarlens}[title={What you've built}]
One more word for this chapter.
\end{grammarlens}

\subsection*{Your turn}
\begin{itemize}
\raggedright
  \item \emph{Say it:} \emph{golódnyy}
  \item \emph{Say it:} \emph{golódnyy}, once more
  \item \emph{Say it:} say \emph{golódnyy}
  \item \emph{Recall:} say the Russian for alive, then the Russian for dead, then say \emph{golódnyy} again
\end{itemize}

\begin{tcolorbox}[breakable,colback=teal!4,colframe=teal!35!black,title={Before you move on}]
What does \ru{голодный} mean? (\textquotedblleft{}Hungry\textquotedblright{}.) Say it once more.
\end{tcolorbox}
\section[\texorpdfstring{\ru{сытый} (sýtyy)}{сытый (sýtyy)}]{\ru{сытый} (sýtyy) — well-fed, full up}

\label{lesson:RU-C158-sytyy}



\begin{quote}
Before the new one: say the Russian for dead, then the Russian for hungry.
\end{quote}

\subsection*{You'll want to know: \ru{сытый}}
\textbf{\ru{сытый}} (\emph{sýtyy}) — \textquotedblleft{}well-fed, full up\textquotedblright{}.

\begin{grammarlens}[title={What you've built}]
One more word for this chapter.
\end{grammarlens}

\subsection*{Your turn}
\begin{itemize}
\raggedright
  \item \emph{Say it:} \emph{sýtyy}
  \item \emph{Say it:} \emph{sýtyy}, once more
  \item \emph{Say it:} say \emph{sýtyy}
  \item \emph{Recall:} say the Russian for dead, then the Russian for hungry, then say \emph{sýtyy} again
\end{itemize}

\begin{tcolorbox}[breakable,colback=teal!4,colframe=teal!35!black,title={Before you move on}]
What does \ru{сытый} mean? (\textquotedblleft{}Well-fed, full up\textquotedblright{}.) Say it once more.
\end{tcolorbox}
\section[\texorpdfstring{\ru{знакомый} (znakómyy)}{знакомый (znakómyy)}]{\ru{знакомый} (znakómyy) — familiar; an acquaintance}

\label{lesson:RU-C158-znakomyy}



\begin{quote}
Before the new one: say the Russian for hungry, then the Russian for well-fed.
\end{quote}

\subsection*{You'll want to know: \ru{знакомый}}
\textbf{\ru{знакомый}} (\emph{znakómyy}) — \textquotedblleft{}familiar; an acquaintance\textquotedblright{}.

\begin{grammarlens}[title={What you've built}]
One more word for this chapter.
\end{grammarlens}

\subsection*{Your turn}
\begin{itemize}
\raggedright
  \item \emph{Say it:} \emph{znakómyy}
  \item \emph{Say it:} \emph{znakómyy}, once more
  \item \emph{Say it:} say \emph{znakómyy}
  \item \emph{Recall:} say the Russian for hungry, then the Russian for well-fed, then say \emph{znakómyy} again
\end{itemize}

\begin{tcolorbox}[breakable,colback=teal!4,colframe=teal!35!black,title={Before you move on}]
What does \ru{знакомый} mean? (\textquotedblleft{}Familiar; an acquaintance\textquotedblright{}.) Say it once more.
\end{tcolorbox}
